\chapter{Supervision and Guidance Strategies for SVG Generation}
\label{ch:supervision}

Representation defines how vector graphics are encoded during generation. Supervision determines what happens to those representations during training or optimization. In SVG generation, the supervision signal is whatever tells the model whether its current output is getting closer to the goal or further away. That signal might come from a target image, a text prompt scored by CLIP, a diffusion model's learned prior, or the structure of a training dataset itself.

What makes supervision in this domain unusual is that it interacts with representation. The same loss function applied to explicit B\'{e}zier paths and to implicit neural fields produces different geometric outcomes, because gradients propagate differently through each representation. This is why Chapters~\ref{ch:representation} and~\ref{ch:supervision} need to be read together: neither dimension explains the structural properties of the output on its own.

Table~\ref{tab:supervision_overview} provides a complete overview of which signal each model uses and when it is applied.


\section{Reconstruction-Based Supervision}
\label{sec:recon_supervision}

Render the SVG, compare it to a target image, update the parameters. That is the entire loop. Reconstruction-based supervision is the most direct application of the losses described in Section~\ref{sec:recon_losses}, and it works through DiffVG~\cite{liDifferentiableVectorGraphics2020a}.

Because the target is known, gradients provide clear information about what should change. When datasets contain clean SVG examples, models trained this way can absorb typical path counts and grouping patterns from the data. DeepSVG, DeepIcon, and Im2Vec all use reconstruction objectives~\cite{carlierDeepSVGHierarchicalGenerative2020, bingDeepIconHierarchicalNetwork2024, reddyIm2VecSynthesizingVector2021}, each with a different combination of losses (Table~\ref{tab:supervision_overview}). Once trained, generation is a single forward pass with no further optimization (Chapter~\ref{ch:Landscape}).

The structural quality of the output comes from the dataset, not from the loss. The loss only measures visual similarity. A model trained on messy SVGs with redundant paths would learn to reproduce that mess just as faithfully. A model can achieve low reconstruction error while producing weak hierarchy or overlapping geometry, because pixel-level matching does not care about what the SVG file looks like on the inside. Reconstruction supervision also requires a reference image, so it cannot support text-to-SVG generation on its own.


\section{Semantic Guidance with Vision--Language Models}
\label{sec:clip_guidance}

When no reference image exists, CLIP-based guidance (Section~\ref{sec:clip}) can score how well a rendered SVG matches a text prompt. This makes open-ended generation possible without paired SVG datasets.

The cost is geometric. CLIP measures alignment in embedding space, not structural correctness. Optimization can introduce extra primitives or distort geometry to push the embedding score higher, and there is nothing in the signal that penalizes fragmented or uneditable output. In practice, CLIP guidance is rarely used alone. It typically serves as an early-stage alignment signal before diffusion-based supervision takes over, as in VectorFusion~\cite{jainVectorFusionTexttoSVGAbstracting2022} and early SVGDreamer configurations. CLIP gets you in the right semantic neighborhood. Getting the geometry right requires a stronger signal.


\section{Diffusion-Based Supervision}
\label{sec:diffusion_supervision}

Diffusion-based supervision has become the dominant strategy for text-to-SVG generation. Instead of comparing to a reference image or scoring embedding similarity, these methods use a pretrained diffusion model as a perceptual prior. The diffusion model has already learned what images should look like. The SVG parameters are updated to produce renderings that the diffusion model considers plausible, conditioned on the text prompt.

The mechanism is Score Distillation Sampling (SDS), with Variational Score Distillation (VSD) introduced by SVGDreamer~\cite{xingSVGDreamerTextGuided2024} to address the oversaturation that standard SDS tends to produce. Both are described in Section~\ref{sec:sds_vsd}. What matters for this chapter is what each variant does to the geometry.

VectorFusion and NiVeL use standard SDS~\cite{jainVectorFusionTexttoSVGAbstracting2022, thamizharasanNIVeLNeuralImplicit2024}. NeuralSVG adds a fine-tuned LoRA adapter on top of SDS~\cite{polaczekNeuralSVGImplicitRepresentation2025}. SVGDreamer replaces SDS entirely with VPSD~\cite{xingSVGDreamerTextGuided2024}. Figure~\ref{fig:supervision_comparison} shows the outputs side by side from the same prompt. VectorFusion produces over-smoothed geometry with 64 shapes. SVGDreamer improves diversity but uses 256 shapes to get there. NeuralSVG reaches comparable visual quality with 16 shapes. The supervision variant matters, but so does the representation it operates through.

Diffusion gradients operate in pixel space and reward perceptual realism. They do not reward geometric simplicity. Without structural constraints, optimization tends to add more primitives to satisfy visual objectives. Inference cost is also high: every generation requires repeated rendering and diffusion evaluation cycles, unlike dataset-trained models that produce output in a single pass.

\begin{figure}[ht]
    \centering
    \begin{tabular}{ccc}
    \textbf{VectorFusion} & \textbf{SVGDreamer} & \textbf{NeuralSVG} \\
    \small{SDS} & \small{VPSD} & 
    \small{SDS + LoRA} \\[4pt]
    
\includegraphics[width=0.3\textwidth]{vectorfusion_img.png} &
    
\includegraphics[width=0.3\textwidth]{svgdreamer256_img.png} &
    
\includegraphics[width=0.3\textwidth]{neuralsvg16_img.png} \\
    \end{tabular}
    \caption{Three diffusion-based supervision variants generating 
    from the same prompt. VectorFusion (SDS, 64 shapes) produces 
    over-smoothed geometry. SVGDreamer (VPSD, 256 shapes) improves 
    diversity but uses four times as many primitives. NeuralSVG 
    (SDS with LoRA, 16 shapes) achieves comparable visual quality 
    with the fewest shapes. Images reproduced 
    from~\cite{xingSVGFusionScalableTexttoSVG2025, 
    polaczekNeuralSVGImplicitRepresentation2025}.}
    \label{fig:supervision_comparison}
\end{figure}


\section{Dataset Supervision vs.\ Per-Instance Supervision}
\label{sec:dataset_vs_instance}

When supervision is applied changes what kind of output you get.

Dataset-supervised models learn from collections of SVG graphics during training. DeepSVG~\cite{carlierDeepSVGHierarchicalGenerative2020}, DeepIcon~\cite{bingDeepIconHierarchicalNetwork2024}, and Im2Vec~\cite{reddyIm2VecSynthesizingVector2021} are trained on reconstruction objectives against SVG data, and once trained, generation requires no further optimization. LayerTracer~\cite{songLayerTracerCognitiveAlignedLayered2025} takes this further: it trains on a synthetically generated serpentine dataset where vector graphics are constructed in a deliberate layer order that mimics how designers actually build graphics (Figure~\ref{fig:serpentine}). The dataset structure itself is the supervision prior. The model learns to produce layered outputs not because a loss function rewards layering, but because every training example was built that way.

\begin{figure}[ht]
    \centering
    \includegraphics[width=0.85\textwidth]{latertracer trng.png}
    \caption{Step-by-step construction sequence from the Serpentine 
    dataset used to train 
    LayerTracer~\cite{songLayerTracerCognitiveAlignedLayered2025}. 
    Each stage adds a new layer, mimicking how a designer would 
    build the graphic incrementally. The dataset structure itself 
    encodes the supervision prior.}
    \label{fig:serpentine}
\end{figure}

Per-instance supervision works on a completely different timescale. LIVE~\cite{maLayerwiseImageVectorization2022}, SAMVG~\cite{zhuSAMVGMultistageImage2023}, VectorFusion~\cite{jainVectorFusionTexttoSVGAbstracting2022}, SVGDreamer~\cite{xingSVGDreamerTextGuided2024}, NeuralSVG~\cite{polaczekNeuralSVGImplicitRepresentation2025}, NiVeL~\cite{thamizharasanNIVeLNeuralImplicit2024}, and T2V-NPR~\cite{zhangTexttoVectorGenerationNeural2024} all optimize vector parameters from scratch for each input. No reusable model is trained. Structure emerges dynamically in response to each prompt or image, which is flexible but computationally expensive and introduces variability from initialization choices.

SVGFusion~\cite{xingSVGFusionScalableTexttoSVG2025} and T2V-NPR~\cite{zhangTexttoVectorGenerationNeural2024} combine both: they learn latent representations from datasets during training, then apply optimization-based refinement at inference. The structural consistency comes from the training data. The flexibility comes from per-instance optimization. The complexity comes from both.


\begin{table}[ht]
\centering
\caption{Supervision signals used by each of the fourteen models 
analyzed in this thesis.}
\label{tab:supervision_overview}
\begin{tabular}{l l l}
\hline
\textbf{Model} & \textbf{Supervision Signal} & \textbf{Applied} \\
\hline
LIVE & L1 + perceptual (VGG) + Xing loss & Inference \\
SAMVG & L1 + IoU + SAM segmentation init & Inference \\
Im2Vec & L2 via differentiable rendering & Training \\
VectorFusion & SDS from Stable Diffusion & Inference \\
SVGDreamer & VPSD from Stable Diffusion & Inference \\
SVGDreamer++ & VPSD + HIVE masks + reward model & Inference \\
NeuralSVG & SDS + fine-tuned LoRA adapter & Inference \\
NiVeL & SDS from Stable Diffusion & Inference \\
DeepSVG & Cross-entropy + L1 + KL divergence & Training \\
DeepIcon & Cross-entropy + CLIP embedding & Training \\
SVGFusion & KL + latent-space denoising (VS-DiT) & Training \\
T2V-NPR & VSD + neural path VAE & Both \\
LayerTracer & Denoising + serpentine construction data & Training \\
\hline
\end{tabular}
\end{table}


\section{Initialization and Structural Priors}
\label{sec:init_priors}

Not everything that shapes a generated SVG comes from a loss function. Where optimization starts matters at least as much as what it optimizes for.

SAMVG~\cite{zhuSAMVGMultistageImage2023} is the clearest example. It uses the Segment Anything Model to define semantically meaningful regions before any vector optimization begins. The loss function is still pixel reconstruction, same as LIVE. But SAMVG's outputs have cleaner boundaries and less path fragmentation, because the starting point already encodes where objects are. The supervision signal did not change. The initialization did, and that made the structural difference. Not without cost though: SAMVG's outputs are structurally cleaner but visually less faithful to the target (Figure~\ref{fig:init_comparison}). Even initialization feeds into the realism-editability trade-off.

\begin{figure}[ht]
    \centering
    \begin{tabular}{cc}
    \textbf{LIVE} & \textbf{SAMVG} \\
    \textbf{Random init} & \textbf{SAM segmentation init} \\[4pt]
    \includegraphics[width=0.4\textwidth]{live_img.png} &
    \includegraphics[width=0.4\textwidth]{sam_img.png} \\
    \end{tabular}
    \caption{Both models use pixel reconstruction loss via 
    differentiable rasterization. LIVE initializes randomly and 
    produces a visually closer match to the target through dense 
    overlapping paths, but the resulting file is fragmented and 
    difficult to edit. SAMVG initializes from segmentation masks 
    and produces cleaner object boundaries with fewer paths, at 
    the cost of reduced visual 
    fidelity~\cite{zhuSAMVGMultistageImage2023}.}
    \label{fig:init_comparison}
\end{figure}

SVGFusion and T2V-NPR both start optimization from latent embeddings shaped by prior training, so the optimizer never begins from random noise. LayerTracer goes further: its training data encodes construction order (Figure~\ref{fig:serpentine}), so the model learns to produce layered outputs without any explicit structural loss. The dataset itself is the structural prior.

This explains something that supervision type alone cannot: why models using the same diffusion-based loss produce very different structural outcomes. When meaningful priors are absent, diffusion optimization piles on primitives. When structured initialization is present, the same loss produces cleaner geometry with fewer paths. Supervision tells you what the model optimizes for. Initialization tells you where it starts. That starting point often determines whether the result is editable.


\section{Summary}
\label{sec:supervision_summary}

The recurring pattern across all supervision strategies is the same: the signal prioritizes visual objectives while structural properties emerge indirectly, if they emerge at all. Initialization can partly compensate (Section~\ref{sec:init_priors}), but no supervision signal in the current landscape directly rewards editability.
